\chapter{年度不明の問題}
\label{ch:unknown}

\begin{memo*}
元資料に年度表記がないため，推測で年度を付けずに掲載する。写真の問題番号から，少なくとも二つの試験冊子が混在している。
\end{memo*}

\section{専門基礎科目と思われる問題}

\begin{problem}{行列｜直交対角化}{unknown-1}
行列
\[
P=\begin{pmatrix}
1/\sqrt3&1/\sqrt2&p\\
q&r&s\\
-1/\sqrt3&1/\sqrt2&-1/\sqrt6
\end{pmatrix}
\qquad(p,q,r,s\in\R)
\]
は直交行列であるとする。
\begin{enumerate}
  \item[(1)] $p,q,r,s$ の取りうる組み合わせを全て求めよ。
  \item[(2)] 行列 ${}^{\mathsf T}(P^{10})\,P^{19}\,({}^{\mathsf T}P)^{10}$ を求めよ。
  \item[(3)] 行列 $A$ は $P$ で対角化されるとし，$A$ の固有多項式 $g_A(x)$ は $x^3+\sqrt[10]{2}\,x^2$ であるとする。また，ベクトル $(1,0,1)^{\mathsf T}$ は固有値 $-\sqrt[10]{2}$ に属する $A$ の固有ベクトルであるとする。このとき $A^{10}$ を求めよ。
  \item[(4)] 直交行列で対角化される行列は対称行列であることを証明せよ。
\end{enumerate}
\end{problem}
\begin{memo*}
\textbf{解答の見通し。}
直交行列は $P^{\mathsf T}P=I$ を満たす行列である。これは、各列の長さが1で、
異なる二列の内積が0であることと同値なので、その条件を式にして $p,q,r,s$ を順に決める。
$P^{\mathsf T}=P^{-1}$ だから、中央の行列積は同じ底 $P$ の整数乗へ直して指数を加減できる。

$A$ の固有値は固有多項式から0,0,$-\sqrt[10]{2}$ と分かる。
非零固有値の単位固有ベクトルを $u$ とし、残り二本の単位固有ベクトルを加えて直交行列 $Q$ を作ると、
$A=Q\diag(-\sqrt[10]{2},0,0)Q^{\mathsf T}$ である。この積を列ごとに展開すると
$A=-\sqrt[10]{2}\,uu^{\mathsf T}$ となり、10乗を計算できる。
\end{memo*}
\begin{proof}[解答]
\begin{enumerate}
  \item[(1)] 列ベクトルの正規直交条件から
  \[
  r=0,\qquad q=\frac{\varepsilon}{\sqrt3},\qquad
  p=\frac1{\sqrt6},\qquad
  s=-\varepsilon\sqrt{\frac23}
  \qquad(\varepsilon=\pm1).
  \]
  これらはいずれも第3列の長さも1にするので、全ての組である。

  \item[(2)] $P^{\mathsf T}=P^{-1}$ より
  \[
  (P^{10})^{\mathsf T}P^{19}(P^{\mathsf T})^{10}
  =P^{-10}P^{19}P^{-10}=P^{-1}=P^{\mathsf T}.
  \]

  \item[(3)] 固有値は $0,0,-\sqrt[10]{2}$ である。
  \[
  u=\frac1{\sqrt2}(1,0,1)^{\mathsf T}
  \]
  とすると、直交対角化より $A=-\sqrt[10]{2}\,uu^{\mathsf T}$。
  $u^{\mathsf T}u=1$ だから $(uu^{\mathsf T})^k=uu^{\mathsf T}$ であり、
  \[
  A^{10}=2uu^{\mathsf T}
  =\begin{pmatrix}1&0&1\\0&0&0\\1&0&1\end{pmatrix}.
  \]

  \item[(4)] 直交行列 $Q$ と実対角行列 $D$ により
  $Q^{-1}AQ=D$ と書けるなら $A=QDQ^{\mathsf T}$。従って
  \[
  A^{\mathsf T}=QD^{\mathsf T}Q^{\mathsf T}=QDQ^{\mathsf T}=A.
  \]
\end{enumerate}
\end{proof}

\begin{problem}{微分積分｜積分記号下の微分}{unknown-2}
$a,b,c,d\in\R$ に対し $a<b$, $c<d$ を仮定する。$D$ は $\R^2$ の開集合で $[a,b]\times[c,d]$ を含むとする。$f(x,y)$ を $D$ 上の $C^1$ 級関数とする。
\begin{enumerate}
  \item[(1)] $[a,b]$ 上で
  \[
  \frac{\d}{\d x}\int_c^d f(x,y)\,\d y=\int_c^d\frac{\partial f}{\partial x}(x,y)\,\d y
  \]
  が成り立つことを示せ。ただし $G(x)=\int_c^d\frac{\partial f}{\partial x}(x,y)\,\d y$ が $[a,b]$ 上で連続であることを仮定して良い。
  \item[(2)] $\phi(x),\psi(x)$ は $[a,b]$ を含む開区間上の $C^1$ 級関数で，任意の $x\in[a,b]$ に対し $\phi(x),\psi(x)\in[c,d]$ が成り立つとする。このとき
  \[
  \frac{\d}{\d x}\int_{\phi(x)}^{\psi(x)}f(x,y)\,\d y
  =f(x,\psi(x))\frac{\d\psi}{\d x}(x)-f(x,\phi(x))\frac{\d\phi}{\d x}(x)
  +\int_{\phi(x)}^{\psi(x)}\frac{\partial f}{\partial x}(x,y)\,\d y
  \]
  が成り立つことを示せ。
  \item[(3)] $x>0$ に対し
  \[
  p(x)=\int_{x^2}^{x^3}\frac{\sin(xy)}y\,\d y
  \]
  とおくとき，$\dfrac{\d p}{\d x}(x)$ を求めよ。
\end{enumerate}
\end{problem}
\begin{memo*}
\textbf{解答の見通し。}
固定端点とは積分区間 $[c,d]$ が $x$ によらない場合である。
差商 $\{F(x+h)-F(x)\}/h$ を書き、$f$ の差だけを積分した形へ直す。
一変数の微積分学の基本定理を $x$ 方向へ使うと、その差商は $\partial f/\partial x$ の平均になる。
$\partial f/\partial x$ の連続性により平均が一様に $\dfrac{\partial f}{\partial x}(x,y)$ へ近づくため、極限と積分を交換できる。

可変端点では $Q(x,v)=\int_c^v f(x,y)\,\d y$ と置き、
$Q(x,\psi(x))-Q(x,\phi(x))$ を二変数の連鎖律で微分する。
これにより、上端の移動、下端の移動、被積分関数自身の変化という三つの項が現れる。
\end{memo*}
\begin{proof}[解答]
\begin{enumerate}
  \item[(1)] $F(x)=\int_c^d f(x,y)\,\d y$ とおく。差商は
  \[
  \frac{F(x+h)-F(x)}h
  =\int_c^d\int_0^1 f_x(x+th,y)\,\d t\d y.
  \]
  $f_x$ の連続性により $h\to0$ で被積分関数は一様に $f_x(x,y)$ へ収束するから
  \[
  F'(x)=\int_c^d f_x(x,y)\,\d y.
  \]

  \item[(2)] $Q(x,v)=\int_c^v f(x,y)\,\d y$ とおくと
  \[
  Q_x(x,v)=\int_c^v f_x(x,y)\,\d y,\qquad Q_v(x,v)=f(x,v).
  \]
  $\int_{\phi(x)}^{\psi(x)}f(x,y)\,\d y
  =Q(x,\psi(x))-Q(x,\phi(x))$ に連鎖律を用いれば
  \[
  \frac{\d}{\d x}\int_{\phi(x)}^{\psi(x)}f(x,y)\,\d y
  =f(x,\psi(x))\psi'(x)-f(x,\phi(x))\phi'(x)
  +\int_{\phi(x)}^{\psi(x)}f_x(x,y)\,\d y.
  \]

  \item[(3)] $f(x,y)=\sin(xy)/y$, $\phi(x)=x^2$, $\psi(x)=x^3$ として (2) を使うと
  \begin{align*}
  p'(x)
  &=\frac{3\sin(x^4)-2\sin(x^3)}x
    +\int_{x^2}^{x^3}\cos(xy)\,\d y\\
  &=\frac{4\sin(x^4)-3\sin(x^3)}x.
  \end{align*}
\end{enumerate}
\end{proof}

\begin{problem}{距離空間｜コンパクト集合の近傍}{unknown-3}
$(X,d)$ を距離空間とし，$K$ を $X$ のコンパクトな部分集合とする。関数 $\rho:X\to\R$ を
\[
\rho(x)=\inf\{d(x,z)\mid z\in K\}\qquad(x\in X)
\]
で定める。
\begin{enumerate}
  \item[(1)] すべての $x,y\in X$ に対して $|\rho(x)-\rho(y)|\le d(x,y)$ が成り立つことを示せ。
  \item[(2)] 任意の $x\in X$ に対して，ある $a\in K$ が存在して $\rho(x)=d(x,a)$ を満たすことを示せ。
  \item[(3)] $L$ は $K$ を含む $X$ の開集合とする。このとき $K_r=\{x\in X\mid\rho(x)<r\}\subset L$ となる実数 $r>0$ が存在することを示せ。
\end{enumerate}
\end{problem}
\begin{memo*}
\textbf{解答の見通し。}
$\rho(x)$ は、$x$ と $K$ の各点との距離の下限である。
三角不等式を全ての $z\in K$ について書き、下限を取ると
$|\rho(x)-\rho(y)|\le d(x,y)$ が得られる。
最短点の存在は、連続関数 $z\mapsto d(x,z)$ がコンパクト集合 $K$ 上で最小値を取ることから示す。

$K$ を含む開集合 $L$ の外側を $F=X\setminus L$ とする。
$F$ が空なら $L=X$ なので、どの $r>0$ を選んでも $\{\rho<r\}\subset L$ である。
$F$ が空でなければ、各 $z\in K$ と閉集合 $F$ の距離は正であり、
それを $K$ 上で最小化した値 $\delta$ も正になる。$0<r<\delta$ を選べば、
$K$ から距離 $r$ 未満の点が $F$ に入ることはできない。
\end{memo*}
\begin{proof}[解答]
\begin{enumerate}
  \item[(1)] 任意の $z\in K$ に対し
  \[
  \rho(x)\le d(x,z)\le d(x,y)+d(y,z).
  \]
  $z$ について下限を取り、$x,y$ を交換した式と合わせれば
  \[
  |\rho(x)-\rho(y)|\le d(x,y).
  \]

  \item[(2)] $z\mapsto d(x,z)$ は $K$ 上連続である。$K$ はコンパクトだから
  最小値を取る $a\in K$ があり、
  \[
  \rho(x)=\min_{z\in K}d(x,z)=d(x,a).
  \]

  \item[(3)] 各 $z\in K$ に対し、$L$ が開なので
  $B(z,2\varepsilon_z)\subset L$ となる $\varepsilon_z>0$ を取る。
  コンパクト性より
  \[
  K\subset B(z_1,\varepsilon_{z_1})\cup\cdots\cup
  B(z_m,\varepsilon_{z_m}).
  \]
  $r=\min_i\varepsilon_{z_i}>0$ とする。$\rho(x)<r$ なら (2) より
  $d(x,a)<r$ となる $a\in K$ がある。ある $i$ で
  $d(a,z_i)<\varepsilon_{z_i}$ だから
  \[
  d(x,z_i)<r+\varepsilon_{z_i}\le2\varepsilon_{z_i}.
  \]
  従って $x\in L$、すなわち $K_r\subset L$。
\end{enumerate}
\end{proof}

\section{専門科目と思われる問題}

\begin{problem}{回転面｜トーラス}{unknown-4}
$R>r>0$ とする。$\R^3$ において $xz$ 平面上の円
\[
\{(x,0,z)\in\R^3\mid(x-R)^2+z^2=r^2\}
\]
を $z$ 軸のまわりに回転して得られる曲面を $S$ とする。
\begin{enumerate}
  \item[(1)] 曲面 $S$ のガウス曲率 $K$ と平均曲率 $H$ を求めよ。
  \item[(2)] 曲面 $S$ 上の測地線の例を一つあげ，それが測地線となることを説明せよ。
\end{enumerate}
\end{problem}
\begin{memo*}
\textbf{解答の見通し。}
回転角 $u$ と母円の角 $v$ で標準的にパラメータ表示する。
二つの座標ベクトルは直交するため、第1・第2基本形式から主曲率を直接読める。
測地線の例には外側の赤道 $v=0$ を選び、
その加速度が曲面の法線方向だけであることを確認する。
\end{memo*}
\begin{proof}[解答]
\begin{enumerate}
\item[(1)]
\[
X(u,v)=((R+r\cos v)\cos u,(R+r\cos v)\sin u,r\sin v)
\]
とおく。まず偏微分は
\[
X_u=(-(R+r\cos v)\sin u,(R+r\cos v)\cos u,0),
\]
\[
X_v=(-r\sin v\cos u,-r\sin v\sin u,r\cos v).
\]
従って第1基本形式の係数は
\[
E=(R+r\cos v)^2,\qquad F=0,\qquad G=r^2.
\]
$R>r$ より $R+r\cos v>0$ なので $EG-F^2>0$、従ってこれは正則な曲面である。
外向き単位法線を
\[
N=(\cos v\cos u,\cos v\sin u,\sin v)
\]
と取る。二階偏微分と内積を取ると
\[
L=X_{uu}\cdot N=-(R+r\cos v)\cos v,\quad
M=X_{uv}\cdot N=0,\quad
N_2=X_{vv}\cdot N=-r.
\]
$F=M=0$ なので座標方向が主方向であり、主曲率は
\[
k_1=\frac LE=-\frac{\cos v}{R+r\cos v},\qquad
k_2=\frac{N_2}G=-\frac1r.
\]
従って
\[
K=\frac{\cos v}{r(R+r\cos v)},\qquad
H=-\frac{R+2r\cos v}{2r(R+r\cos v)}.
\]
\item[(2)] $v=0$ と固定した外側の赤道を $\gamma(u)=X(u,0)$ とおくと
\[
\gamma''(u)=-(R+r)(\cos u,\sin u,0)=-(R+r)N(u,0).
\]
加速度の接平面成分が0、すなわち測地的曲率が0なので、この円は測地線である。
\end{enumerate}
\end{proof}

\begin{problem}{Euler 型微分方程式}{unknown-5}
$p,q\in\R$ とする。微分方程式に対する初期値問題
\[
x^2\frac{\d^2y}{\d x^2}+px\frac{\d y}{\d x}+qy=0\quad(x\ge1),
\qquad
y(1)=0,\quad \frac{\d y}{\d x}(1)=1
\tag{1}
\]
を考える。
\begin{enumerate}
  \item[(1)] $p\ne1$ かつ $q=0$ とするとき，(1) の解 $y(x)$ を $p$ を用いて表せ。
  \item[(2)] 変数変換 $x=e^t$ によって，(1) の第一式を変形して得られる方程式を求めよ。
  \item[(3)] $p=2$ とするとき，(1) の解 $y(x)$ を $q$ を用いて表せ。
  \item[(4)] (3) で得られた解 $y(x)$ について，$y(x)=0$ となる $x>1$ が存在するような $q$ の範囲を求め，そのような $x>1$ の最小値を $q$ を用いて表せ。
\end{enumerate}
\end{problem}
\begin{memo*}
\textbf{解答の見通し。}
Euler 型方程式は対数変数 $t=\log x$ へ移すと定数係数方程式になる。
$p=2$ では判別式 $d=1-4q$ の符号で、双曲線関数・重根・三角関数の三場合に分ける。
$x>1$ は $t>0$ に対応し、零点が繰り返し現れるのは三角関数になる
$d<0$、すなわち $q>1/4$ の場合だけである。
\end{memo*}
\begin{proof}[解答]
\begin{enumerate}
\item[(1)] $q=0$ では $x^2y''+pxy'=0$ である。$z=y'$ とおくと
\[
xz'+pz=0,\qquad \frac{z'}z=-\frac px.
\]
従って $z=Cx^{-p}$。初期条件 $y'(1)=1$ から $C=1$ で、$y(1)=0$ を使って
\[
y(x)=\int_1^x s^{-p}\,\d s=\frac{x^{1-p}-1}{1-p}\qquad(p\ne1).
\]

\item[(2)] $x=e^t$、$Y(t)=y(e^t)$ とおく。連鎖律から
\[
Y'(t)=e^ty'(e^t)=xy'(x),
\]
\[
Y''(t)=xy'(x)+x^2y''(x),
\]
従って $x^2y''=Y''-Y'$ である。元の方程式は
\[
Y''+(p-1)Y'+qY=0,\qquad Y(0)=0, Y'(0)=1.
\]
\item[(3)] $p=2$ では特性方程式が $\lambda^2+\lambda+q=0$ となる。
$d=1-4q$ とおく。

$d>0$ のとき、特性根は $(-1\pm\sqrt d)/2$ なので
\[
Y(t)=e^{-t/2}\left(C_1e^{\sqrt d\,t/2}+C_2e^{-\sqrt d\,t/2}\right).
\]
$Y(0)=0$ から $C_2=-C_1$。従って
$Y(t)=2C_1e^{-t/2}\sinh(\sqrt d\,t/2)$ であり、
\[
Y'(0)=2C_1\cdot\frac{\sqrt d}{2}=C_1\sqrt d=1
\]
から $C_1=1/\sqrt d$ を得る。

$d=0$ のときは重根 $-1/2$ を持つので
$Y(t)=(C_1+C_2t)e^{-t/2}$。初期条件から $C_1=0$, $C_2=1$ である。

$d<0$ のとき、$b=\sqrt{-d}$ とおけば特性根は $(-1\pm ib)/2$ であり
\[
Y(t)=e^{-t/2}\left(C_1\cos\frac{bt}{2}+C_2\sin\frac{bt}{2}\right).
\]
$Y(0)=0$ から $C_1=0$、さらに
$Y'(0)=C_2b/2=1$ から $C_2=2/b$ となる。

最後に $t=\log x$、$e^{-t/2}=x^{-1/2}$ を戻すと
\[
y=x^{-1/2}\begin{cases}
\dfrac{2}{\sqrt d}\sinh(\frac{\sqrt d}{2}\log x),&d>0,\\
\log x,&d=0,\\
\dfrac{2}{\sqrt{-d}}\sin(\frac{\sqrt{-d}}2\log x),&d<0.
\end{cases}
\]
\item[(4)] $d\ge0$ の場合、$x>1$ では $\log x>0$ であり、$\sinh u$ と $u$ は正なので零点を持たない。
$d<0$、すなわち $q>1/4$ の場合は
\[
\sin\left(\frac{\sqrt{4q-1}}2\log x\right)=0
\]
より、正の零点は
\[
\frac{\sqrt{4q-1}}2\log x=k\pi\qquad(k=1,2,\ldots).
\]
従って $x>1$ の零点があるのは $q>1/4$ のときに限り、最小の零点は
$\exp(2\pi/\sqrt{4q-1})$ である。
\end{enumerate}
\end{proof}

\begin{problem}{解析｜弱い特異核}{unknown-6}
$\alpha$ を $0<\alpha<1$ を満たす実数とする。$[0,\infty)$ 上で有界なルベーグ可測関数 $\phi(t)$ に対し，$[0,\infty)$ 上の関数 $f(x)$ を
\[
f(x)=\int_0^x\phi(t)\,\d t
\]
で定める。
\begin{enumerate}
  \item[(1)] $f(x)$ は $[0,\infty)$ 上で連続であることを示せ。
  \item[(2)] $(0,\infty)$ 上の関数 $g(x)$ を
  \[
  g(x)=\int_0^x\frac{f(s)}{(x-s)^\alpha}\,\d s
  \]
  で定める。任意の $x\in(0,\infty)$ に対して
  \[
  g(x)=\frac1{1-\alpha}\int_0^x\phi(t)(x-t)^{1-\alpha}\,\d t
  \]
  が成り立つことを示せ。
  \item[(3)] $g(x)$ を (2) で定めた $(0,\infty)$ 上の関数とする。$\{h_n\}$ は正の実数列で，$n\to\infty$ のとき $h_n\to0$ とする。任意の $x\in(0,\infty)$ に対して
  \[
  \lim_{n\to\infty}\frac{g(x+h_n)-g(x)}{h_n}=\int_0^x\frac{\phi(t)}{(x-t)^\alpha}\,\d t
  \]
  が成り立つことを示せ。
\end{enumerate}
\end{problem}
\begin{memo*}
\textbf{解答の見通し。}
$\phi$ の有界性から原始関数 $f$ は Lipschitz 連続になる。
$f(s)=\int_0^s\phi(t)\,\d t$ を代入すると三角形領域の二重積分になるので、
$0<\alpha<1$ による絶対可積分性を確認して Fubini の定理を使う。
導関数候補 $G$ を先に定め、$\int_0^xG(u)\,\d u=g(x)$ を示して
微積分学の基本定理へ帰着する。
\end{memo*}
\begin{proof}[解答]
\begin{enumerate}
\item[(1)] $|f(y)-f(x)|\le\norm\phi_\infty|y-x|$ なので $f$ は連続である。
\item[(2)]$f(s)=\int_0^s\phi(t)dt$ を代入し，Fubini で積分順序を交換すると
\begin{align*}
g(x)
&=\int_0^x\int_0^s\phi(t)(x-s)^{-\alpha}\,\d t\d s\\
&=\int_0^x\phi(t)\int_t^x(x-s)^{-\alpha}\,\d s\d t.
\end{align*}
ここで $|\phi(t)|(x-s)^{-\alpha}$ は三角形 $0\le t\le s\le x$ 上可積分である。実際、
\[
\int_0^x\int_0^s\norm\phi_\infty(x-s)^{-\alpha}\,\d t\d s
=\norm\phi_\infty\int_0^x s(x-s)^{-\alpha}\,\d s<\infty
\]
なので Fubini の定理を適用できる。内側を $u=x-s$ で積分すると
\[
\int_t^x(x-s)^{-\alpha}\,\d s
=\int_0^{x-t}u^{-\alpha}\,\d u
=\frac{(x-t)^{1-\alpha}}{1-\alpha},
\]
従って指定された表示式を得る。

\item[(3)] 導関数の候補を
\[
G(x)=\int_0^x\phi(t)(x-t)^{-\alpha}\,\d t
\]
とおく。$0<\alpha<1$ なのでこの積分は絶対収束する。さらに Fubini の定理により
\begin{align*}
\int_0^xG(u)\,\d u
&=\int_0^x\int_0^u\phi(t)(u-t)^{-\alpha}\,\d t\d u\\
&=\int_0^x\phi(t)\int_t^x(u-t)^{-\alpha}\,\d u\d t\\
&=\frac1{1-\alpha}\int_0^x\phi(t)(x-t)^{1-\alpha}\,\d t\\
&=g(x).
\end{align*}
$G$ は局所連続である。実際、核 $u^{-\alpha}\mathbf1_{(0,L)}(u)$ は
$L^1(\R)$ に属し、その平行移動は $L^1$ で連続、$\phi$ は有界だからである。
従って微積分学の基本定理を適用でき
\[
g'(x)=G(x)=\int_0^x\phi(t)(x-t)^{-\alpha}\,\d t
\]
を得る。特に問題の極限は $h_n>0$, $h_n\to0$ に対する差商の極限なので、この値に等しい。
\end{enumerate}
\end{proof}
